\chapter{Conclusion}

This thesis investigated the effectiveness of large language models for
Java-to-C++ code translation. Two code-specific model families, StarCoder2 and
Qwen2.5-Coder, were evaluated in both their base and fine-tuned forms using BLEU
score and execution-based evaluation with CodeEditorBench.

The results showed that fine-tuning improved the performance of both models in
terms of textual similarity and functional correctness. This suggests that
task-specific adaptation is an effective way to improve code translation
quality, even when relatively small parameter-efficient fine-tuning methods are
used. The experiments also showed that different model families do not benefit
equally from fine-tuning, indicating that architecture and pre-training strategy
can strongly influence the final translation performance.

The experiments also demonstrated that similarity-based metrics such as BLEU are
not sufficient for evaluating code translation systems on their own. Several
generated programs achieved relatively high BLEU scores while still failing to
compile or execute correctly. This supports earlier research showing that
execution-based evaluation provides a more reliable measure of functional
correctness than token-level similarity metrics.

Although the fine-tuned models performed substantially better than their base
versions, the number of remaining compilation errors and wrong answers indicates
that fully automated code translation is not yet reliable enough for practical
use without human supervision. Current models can already support developers by
generating initial translations, but human review is still necessary to ensure
correctness, maintainability, and compatibility with the target language.

Future work may build upon these findings in several directions. Evaluating
additional programming language pairs could provide a better understanding of
model generalization across different paradigms. Expanding the set of evaluated
models, including larger architectures, may further improve functional
correctness and translation quality. The use of more diverse training datasets
and advanced fine-tuning strategies could enhance robustness. In addition,
combining multiple evaluation metrics and benchmarks would allow for a more
comprehensive assessment of model performance. Finally, investigating
repository-level code translation represents an important step toward more
realistic and practical software migration scenarios.
